\chapter{关于本书与作者} \label{ch:about-author}

这本书能写出来，并不是因为我比谁聪明，而是因为我真的很喜欢做这件事。从小学第一次在机器人兴趣班里看到自己写的代码能控制小车动起来开始，这种喜欢就没有断过。后来学了单片机、做了竞赛、写了这本书，本质上都是同一件事的延续：想亲手把一套系统从头到尾搞明白。

所以这一章不打算写成长流水账，而是挑几件最能解释“这本书为什么是这个风格”的事。如果读者在书里感觉到作者有点“轴”——不喜欢看不懂的封装、宁可自己写也不愿意用现成的库、总把 AI 当苦力使——那都不是性格缺陷，而是从爱好里长出来的习惯。

% ---------------------------------------------------------------------
\section{一路走来的足迹}

我的技术启蒙可以追溯到小学：因为喜欢机器人，自学了 C 语言，也报了机器人兴趣班。中考后的暑假，我第一次接触 Raspberry Pi 3 开发板，尝试自己安装配置 Linux 系统。高中课业紧张，只能在课余学一点 C++ 和 Arduino。现在回头看，那段日子虽然没有参加什么比赛，却帮我打下了 C 系列语言的基础，也让我确定自己以后想走技术这条路。

进入大学后，我从 STC51 单片机开始入门，后来又基于 STC8H 做了智能小车项目，那是我第一次独立完成一个“从硬件到软件都要自己负责”的作品。之后我又参加了全国大学生电子设计竞赛，和不同专业的同学一起守着 MSP430 硬啃，拿到了省二等奖；也参加过 iCAN 的 AI 视觉挑战赛，从 YOLOv11 一路做到全国总决赛。中间还尝试过用国产 RK3588 跑 Linux 做图像识别，当时因为经验不足留下了遗憾，但也正是这段经历，最终催生了我的毕业设计。

这一路走来，我特别感谢每一位教过我的老师、一起熬过夜的队友、始终支持我的家人和朋友。没有他们，我可能早在某一本看不懂的手册面前就放弃了。

% ---------------------------------------------------------------------
\section{毕业设计：给“热爱”造一个“一次性塑料轮子”}

先说说毕业设计。那段时间我面对的是很多做机器人的同学都会遇到的两难：ROS2 的节点通信、EKF 的融合调参、导航精度调优，都需要大量时间去学，而交期只剩很短的时间。硬啃当然可以，但风险是最后连一个能跑的东西都交不出来。

我最初的选择是“正规军”路线：用 ROS2 跑导航，串口协议对接嵌入式底层，同时在边缘端部署视觉模型。听起来每一步都有现成的方案，但真正动手才发现处处是坑——ROS2 的串口通信协议写了半天，收发总是不稳定；嵌入式 C 端要解析 ROS2 发过来的字符指令，光一个健壮的配置文件读取模块就折腾了好久；视觉模型在边缘设备上的推理精度始终达不到预期，调了两个月，白天调参数、晚上改代码，效果依然时好时坏。东西确实在一点一点往前推，但速度太慢了，慢到我心里清楚：按这个节奏，演示视频交付的那天，我手里大概率只有一个“半成品”。

老师当时说了一句话点醒了我：“技术探索的价值不在于一时的输赢，而在于过程中的成长与积淀。但毕业设计的底线是——你得先毕业。”这句话我后来想了很久：不是放弃追求更好的技术，而是要在有限的时间里，先保证自己能交出一份完整的答卷。

于是我做了一个在当时看来“退步”的决定：放弃 ROS2 方案，放弃已经调了两个月的边缘视觉部署路线，换一个更务实的方向。硬件上采用 STM32H743 与 K230 双核异构架构：STM32H743 负责运动控制和数据融合，K230 负责视觉识别，两者之间只用一条高速串口通信。选 K230 的理由很简单——它不是性能最强的，但它部署起来快，官方提供的 SDK 和例程足够把模型跑起来，不需要我再去啃一套复杂的依赖环境。软件上彻底放弃 Linux 和 ROS2，让 AI 帮我生成一个表达式解析器，把复杂的逻辑写成可配置的字符串指令，嵌入式端只需要读文件、解析字符串、执行命令，不再需要维护一套完整的 ROS2 通信栈；同时把 EKF 改写成基于动态打滑度检测的自适应里程计融合算法，用启发式函数判断编码器和 IMU 谁更可信，再做线性加权——没有卡尔曼滤波那样漂亮的理论保证，但每一行融合逻辑我都亲自验证过。

毕设完成后，我甚至把 FreeRTOS 方案也一并放弃了。不是因为 FreeRTOS 不好，而是我发现自己在这个项目里真正需要的是“极简到不可能出错”的裸机循环：一个主循环、几个中断、一条串口收发，出了问题示波器一搭就能定位，不需要去查任务栈溢出、不需要调试优先级反转。裸机方案虽然“原始”，但它让我在最后一个星期里敢改敢跑，因为每一行改动的影响范围我都算得清。

这套方案并不优雅，没有 ROS2 那样漂亮的分布式架构，没有 EKF 那样严谨的数学推导，更谈不上是一个“产品”。它更像是一个用几天时间拼出来的“一次性塑料轮子”——跑完这次演示，它就完成了历史使命。但它有一个无法替代的优点：每一行代码我都看得懂，每一个环节出了问题我都能在两小时内定位，最后它在规定时间内跑通了演示视频。

后来有人问我为什么不直接学 ROS2。我的回答是：ROS2 确实是好东西，我心里一直知道。只是当时手里没有梯子，只能先搭人墙翻过去。等项目结束、有了时间，我还是会回头去啃它——用 AI 帮我分析日志、批量调参，效率一定比当年手工死磕高得多。那次毕设教会我的不是“拒绝正规框架”，而是“在正确的时间做正确的事”：热爱可以等，交期不等。一次性塑料轮子虽然跑不远，但它能让你先过线。过了线，才有资格谈下一次用什么样的轮子跑得更快。

% ---------------------------------------------------------------------
\section{工作以后：宁愿自己写一个“能看透底牌”的轮子}

如果说毕设教会我的是"时间不够时怎么先活下去"，那工作以后遇到的另一件事，教会我的就是——时间够的时候，为什么值得自己造一个重轮子。

当时团队要给产品换UI框架，组长推荐的是基于HelixScreen的方案。那套东西本身不算差，但它本质上只是对LVGL的一层薄封装：没有可视化设计器，每次改界面都像在结构体数组里盲猜坐标——改完编译、下载、看效果，不行再重来。更麻烦的是，这套代码不属于公司，一旦触到深层BUG，你只能钻进别人的封装层里人肉翻源码。

老板今天说"这里挪一下"，明天说"那里改一点"，在这种框架下，每一次修改都是一场探险——你永远不知道这一改会不会触发封装层里某个隐性的副作用。

所以我选择了自研，底层直接挂Qt的QPainter和QPaintDevice绘图层。我不使用QML，但会使用QSS来控制样式；界面设计则直接上Qt Designer做UI设计器——画界面就拖控件，布局用Layout自动撑，生成的就是一份干干净净的.ui文件，本质上是一份XML，所有属性、坐标、层级关系都摊开在文本里，你能直接看到每个控件的x、y、width、height，想改什么直接改属性，完全不用在结构体数组里盲猜。QSS只管样式，不影响UI结构，所以这一层依然没有任何“魔法”：画矩形就是drawRect，写文字就是drawText，UI文件里的每个节点都能追溯到具体属性值。所有行为都在文档和源码里写得明明白白。

这一层没有任何"魔法"：画矩形就是drawRect，写文字就是drawText，UI文件里的每个节点都能追溯到具体属性值。所有行为都在文档和源码里写得明明白白。我用一天理清框架结构，再借助AI Agent帮我填剩下的代码，大约两周把整个框架搭完——组长的原估算是三个月。

表面上看，我用两周的确定性劳动，买断了老板未来无数次"这里改一下"带来的不确定性风险。但更深一层的理由是：我喜欢那种每一层都能看透的感觉。

工具只是手段，真正让我安心的，是我随时知道它为什么这样工作、出问题该去哪里查。UI设计器不是黑盒——它生成的XML我可以直接读懂；Qt绘图层不是黑盒——它的每一个API都有文档和源码支撑。任何我看不透的自动推断，都是埋在项目里的定时炸弹——不是因为它一定会炸，而是因为炸的时候我根本不知道拆哪根线。

% ---------------------------------------------------------------------
\section{写代码时的几个“怪”习惯}

上面两个例子，其实都指向我在书里反复流露出的几个习惯：

\begin{enumerate}
\item 对黑盒零容忍。我承认“哑巴底层工具”，比如 GCC、Qt 绘图层，因为能看透底牌；我拒绝“替你动脑子的智能管家”，因为任何看不透的自动推断，都是埋在项目里的定时炸弹。
\item 长期资产造重轮子，战时只造一次性塑料轮子。核心框架值得花时间自己打磨，换来长久的掌控感；但交期只剩几天时，目标不是优雅，而是活过今晚。
\item AI 是体力劳工，不是军师。我会让 AI 帮我生成代码、填注释、做批量重构，但不会让它替我决定架构。
\end{enumerate}

这些习惯听起来有点固执，但它们的出发点其实很朴素：因为热爱，所以想搞懂；因为想搞懂，所以不接受稀里糊涂地“能用就行”。

% ---------------------------------------------------------------------
\section{感谢与祝福}

书写的过程，也是回望的过程。一路走来，要感谢的人太多太多。

首先要感谢求学路上每一位教导过我的老师，是你们带我走进实验室，让我在动手实践中找到真正的热爱；感谢一起熬夜备赛的队友们，那些互相打气、彼此扶持的夜晚，是我最珍贵的财富。更要深深感谢我的家人，父母用有限的收入为我创造最好的学习环境，哪怕购买开发板和书籍偶尔会带来经济压力，他们也从未有过一句怨言。这份无条件的支持，是我能一直走下去的底气。也感谢每一位同事在日常协作中给予的耐心与帮助，是你们让我从校园走向职场时没有感到太大的落差，反而多了几分从容。

也要感谢翻开这本书的每一位读者。如果书中的某个思路、某段代码能在某个深夜给你一点启发，或者让你觉得“原来这条路真的有人走过”，那便是它存在的全部意义。技术探索从来不是孤军奋战，愿我们都能保持那份最初的好奇与执着，在自己热爱的领域里，一步一个脚印地走下去。

最后，把最真挚的祝福送给所有帮助过我的人，也送给正在阅读这本书的你——愿我们都能造出自己满意的轮子，也愿每一次调试的尽头，都亮起绿灯。
